% GENERATED FILE. Edit canonical lessons, then run npm run generate:books.
% canonical-source-hash: fnv1a64:b61df49b4638de94
% canonical-lessons: KA-C148-kalisu, KA-C148-huduku, KA-C148-kali, KA-C148-kalisu2, KA-C148-jigi

\chapter{To send, To look for, To learn, To teach, To jump}
\label{ch:ka-to-send-to-look-for-to-learn-to-teach-to-jump}
\hlchaptermodality{148}

\begin{chapteropening}
\textbf{By the end of this chapter:} \emph{I can say to send, to look for, to learn, to teach and to jump.}

Complete the last lesson of chapter 148: i can say to send, to look for, to learn, to teach and to jump.
\end{chapteropening}

\section[kaḷisu]{\kn{ಕಳಿಸು} (kaḷisu) — to send}

\label{lesson:KA-C148-kalisu}



\begin{quote}
Before the new one: say the Kannada for to cut, then the Kannada for to break.
\end{quote}

\subsection*{You'll want to know: \kn{ಕಳಿಸು}}
\textbf{\kn{ಕಳಿಸು}} — \emph{kaḷisu} — \textquotedblleft{}to send\textquotedblright{}.

\begin{grammarlens}[title={What you've built}]
One more word for this chapter.
\end{grammarlens}

\subsection*{Your turn}
\begin{itemize}
\raggedright
  \item \emph{Say it:} \emph{kaḷisu}
  \item \emph{Say it:} \emph{kaḷisu}, once more
  \item \emph{Say it:} say \emph{kaḷisu}
  \item \emph{Recall:} say the Kannada for to cut, then the Kannada for to break, then say \emph{kaḷisu} again
\end{itemize}

\begin{tcolorbox}[breakable,colback=teal!4,colframe=teal!35!black,title={Before you move on}]
What does \kn{ಕಳಿಸು} mean? (\textquotedblleft{}To send\textquotedblright{}.) Say it once more.
\end{tcolorbox}
\section[huḍuku]{\kn{ಹುಡುಕು} (huḍuku) — to look for}

\label{lesson:KA-C148-huduku}



\begin{quote}
Before the new one: say the Kannada for to break, then the Kannada for to send.
\end{quote}

\subsection*{You'll want to know: \kn{ಹುಡುಕು}}
\textbf{\kn{ಹುಡುಕು}} — \emph{huḍuku} — \textquotedblleft{}to look for\textquotedblright{}.

\begin{grammarlens}[title={What you've built}]
One more word for this chapter.
\end{grammarlens}

\subsection*{Your turn}
\begin{itemize}
\raggedright
  \item \emph{Say it:} \emph{huḍuku}
  \item \emph{Say it:} \emph{huḍuku}, once more
  \item \emph{Say it:} say \emph{huḍuku}
  \item \emph{Recall:} say the Kannada for to break, then the Kannada for to send, then say \emph{huḍuku} again
\end{itemize}

\begin{tcolorbox}[breakable,colback=teal!4,colframe=teal!35!black,title={Before you move on}]
What does \kn{ಹುಡುಕು} mean? (\textquotedblleft{}To look for\textquotedblright{}.) Say it once more.
\end{tcolorbox}
\section[kali]{\kn{ಕಲಿ} (kali) — to learn}

\label{lesson:KA-C148-kali}



\begin{quote}
Before the new one: say the Kannada for to send, then the Kannada for to look for.
\end{quote}

\subsection*{You'll want to know: \kn{ಕಲಿ}}
\textbf{\kn{ಕಲಿ}} — \emph{kali} — \textquotedblleft{}to learn\textquotedblright{}.

\begin{grammarlens}[title={What you've built}]
One more word for this chapter.
\end{grammarlens}

\subsection*{Your turn}
\begin{itemize}
\raggedright
  \item \emph{Say it:} \emph{kali}
  \item \emph{Say it:} \emph{kali}, once more
  \item \emph{Say it:} say \emph{kali}
  \item \emph{Recall:} say the Kannada for to send, then the Kannada for to look for, then say \emph{kali} again
\end{itemize}

\begin{tcolorbox}[breakable,colback=teal!4,colframe=teal!35!black,title={Before you move on}]
What does \kn{ಕಲಿ} mean? (\textquotedblleft{}To learn\textquotedblright{}.) Say it once more.
\end{tcolorbox}
\section[kalisu]{\kn{ಕಲಿಸು} (kalisu) — to teach}

\label{lesson:KA-C148-kalisu2}



\begin{quote}
Before the new one: say the Kannada for to look for, then the Kannada for to learn.
\end{quote}

\subsection*{You'll want to know: \kn{ಕಲಿಸು}}
\textbf{\kn{ಕಲಿಸು}} — \emph{kalisu} — \textquotedblleft{}to teach\textquotedblright{}.

\begin{grammarlens}[title={What you've built}]
One more word for this chapter.
\end{grammarlens}

\subsection*{Your turn}
\begin{itemize}
\raggedright
  \item \emph{Say it:} \emph{kalisu}
  \item \emph{Say it:} \emph{kalisu}, once more
  \item \emph{Say it:} say \emph{kalisu}
  \item \emph{Recall:} say the Kannada for to look for, then the Kannada for to learn, then say \emph{kalisu} again
\end{itemize}

\begin{tcolorbox}[breakable,colback=teal!4,colframe=teal!35!black,title={Before you move on}]
What does \kn{ಕಲಿಸು} mean? (\textquotedblleft{}To teach\textquotedblright{}.) Say it once more.
\end{tcolorbox}
\section[jigi]{\kn{ಜಿಗಿ} (jigi) — to jump}

\label{lesson:KA-C148-jigi}



\begin{quote}
Before the new one: say the Kannada for to learn, then the Kannada for to teach.
\end{quote}

\subsection*{You'll want to know: \kn{ಜಿಗಿ}}
\textbf{\kn{ಜಿಗಿ}} — \emph{jigi} — \textquotedblleft{}to jump\textquotedblright{}.

\begin{grammarlens}[title={What you've built}]
One more word for this chapter.
\end{grammarlens}

\subsection*{Your turn}
\begin{itemize}
\raggedright
  \item \emph{Say it:} \emph{jigi}
  \item \emph{Say it:} \emph{jigi}, once more
  \item \emph{Say it:} say \emph{jigi}
  \item \emph{Recall:} say the Kannada for to learn, then the Kannada for to teach, then say \emph{jigi} again
\end{itemize}

\begin{tcolorbox}[breakable,colback=teal!4,colframe=teal!35!black,title={Before you move on}]
What does \kn{ಜಿಗಿ} mean? (\textquotedblleft{}To jump\textquotedblright{}.) Say it once more.
\end{tcolorbox}
